\documentclass[11pt]{article}

\usepackage[margin=1in]{geometry}
\usepackage{amsmath,amssymb}
\usepackage{hyperref}
\setlength{\emergencystretch}{3em}

\newcommand{\QU}{\mathbb Q\mathbb U}
\newcommand{\ZU}{\mathbb Z\mathbb U}

\title{Literature Answer To Rosendal's \(\operatorname{Isom}(\QU)\)
Reflexive Fixed-Point Question}
\author{Run \texttt{fa\_banach\_001}, agent \texttt{agent\_lane\_06}}
\date{June 16, 2026}

\begin{document}
\maketitle

\section*{Status}

\textbf{Status: literature already answered.}  This packet records an exact
later literature answer to Problem 66 of Christian Rosendal's
\emph{Large scale geometry of metrisable groups}, arXiv:1403.3106.

\section*{Source Question}

In arXiv:1403.3106, Problem 66 on page 39 asks whether
\(\operatorname{Isom}(\QU)\), the isometry group of the rational Urysohn metric
space, has a continuous affine isometric action on a reflexive Banach space
without a fixed point.

\section*{Answering Theorem}

The later paper Christian Rosendal, \emph{Equivariant geometry of Banach spaces
and topological groups}, arXiv:1610.01070, answers this question negatively.
Corollary 12 on page 5 states the fixed-point conclusion for
\(\operatorname{Isom}(\QU)\).  The detailed proof is Theorem 81 on pages
47--48: if \(E\) is reflexive, or \(E=L^1([0,1])\), and
\[
  G=\operatorname{Isom}(\ZU)\quad\text{or}\quad G=\operatorname{Isom}(\QU),
\]
then every quasi-cocycle \(b:G\to E\) is bounded.  In particular, every affine
isometric action \(G\curvearrowright E\) has a fixed point.

Taking \(G=\operatorname{Isom}(\QU)\) and \(E\) reflexive gives exactly the
answer to Problem 66: no such fixed-point-free continuous affine isometric
action exists.

\section*{Scope}

This packet answers only Problem 66 of arXiv:1403.3106.  The broader Problem
67 from that source, concerning all non-Archimedean Polish groups, is not
settled by this identification.  Also, arXiv:1610.01070 asks a stronger
Question 82 in which the action is affine but not assumed isometric; that
stronger non-isometric variant remains outside this packet.

\section*{Provenance}

The answer is not original to this run.  The run contribution is only the
literature identification: a separate later paper by Rosendal contains the
exact fixed-point theorem answering the earlier extracted question.

\section*{References}

\begin{itemize}
  \item Christian Rosendal, \emph{Large scale geometry of metrisable groups},
  arXiv:1403.3106.
  \item Christian Rosendal, \emph{Equivariant geometry of Banach spaces and
  topological groups}, arXiv:\allowbreak 1610.01070.
\end{itemize}

\end{document}
